\chapter{Request Generation Engine}

\section{Template Composition Model}
\textbf{Decision}: Legal request text is generated using a block-based composition system: a library of small, reusable, prewritten legal-clause blocks (e.g. requester identity statement, legal-basis citation, verification method, deadline/response-time notice), each using Jinja2 for internal variable substitution. A rule set selects and assembles the appropriate ordered blocks for a given request based on its legal framework (Section 9), request type (Section 10), and company-specific requirements, rather than maintaining one monolithic template per framework/request-type/company combination.

\textbf{Rationale}: A single monolithic template per combination of legal framework, request type, and company would combinatorially explode given the specification's near-term expansion targets (GDPR, UK GDPR, CCPA/CPRA, PIPEDA, LGPD, Australia's Privacy Act, see Section 9), and would force largely duplicated legal language to be maintained and reviewed separately in many places. A block-based system lets generic blocks (e.g. requester identity, verification method) be written and reviewed once and reused across every framework and request type that needs them, while framework-specific blocks (e.g. citing the correct statutory article) are added only where genuinely needed. This also makes it easier to reason about completeness: since block selection is rule-driven rather than authored freehand per template, extending support to a new legal framework is primarily a matter of writing its framework-specific blocks and reusing already-reviewed generic ones, rather than re-deriving an entire letter from scratch.
\section{Template Layer Storage}
\textbf{Decision}: Base legal-framework blocks are bundled with the application and versioned with software releases. Company-specific template overrides or additions are stored within the company's YAML record (Section 3.3). User-authored customizations to templates are stored inside the encrypted vault.

\textbf{Rationale}: Each layer is stored alongside the thing that actually owns its lifecycle and review process. Base legal blocks change on the application's own release cadence and warrant legal-accuracy review before shipping, fitting naturally into the signed-release pipeline already established (Section 2.8, Section 26). Company-specific overrides are inherently part of that company's record and benefit from the same community review, verification, and flagging process already defined for company data generally (Section 14), keeping them in the company YAML record ensures a moderator-approved correction to a company's requirements actually reaches template behavior. User customizations are private to that user and have no place outside their own encrypted vault, consistent with the vault's role as the sole store of personal, user-controlled data (Section 2.1).
\section{Request Preview System}
\textbf{Decision}: A rendered preview showing the exact final content of a request, including which specific fields were included as determined by the data minimization engine (Section 7.3), is always shown to the user and is mandatory before a request can be approved. There is no skip or trusted-template bypass.

\textbf{Rationale}: This directly enforces the Draft → Review → User Approval stages of the request lifecycle (Section 11) and the principle that the system never performs sensitive actions without user approval. Given requests are now assembled from composed blocks (Section 9.1) rather than a single fixed template, an always-mandatory preview is what lets the user verify the assembled result is actually correct and complete for their specific case, rather than trusting block selection logic blindly — particularly important the first time a given framework/request-type/company combination is used.


\sentinelchapterend